\section{Persistent and Situated Agent Lifecycle}
\label{sec:agent-lifecycle}

GAWorld represents an agent as a continuing participant in a shared simulation,
not as an isolated chat session.  The representation separates relatively
stable profile material, mutable runtime state, and artifacts accumulated over
simulated time.  Figure~\ref{fig:agent-lifecycle} shows the connected per-tick
control path.

\begin{figure}[t]
  \centering
  \includegraphics[width=\linewidth]{agent_lifecycle.pdf}
  \caption{The configured default agent-step pipeline.  The first and last
  stages are structural boundaries; intermediate stages may be omitted or
  replaced through configuration.  Domain inputs also arrive through direct
  calls and event-bus contribution points, so the diagram is a control-flow
  account rather than a claim that every subsystem is a plugin.}
  \label{fig:agent-lifecycle}
\end{figure}

\subsection{Profile seeds and mutable state}

At construction time, a tabular seed row supplies the identifier,
demographic fields, and initial scalar state, while a Markdown profile supplies
descriptive material such as occupation, personality, daily life, and values.
The current builder merges these sources into a dictionary containing
\texttt{id}, profile fields, \texttt{state}, an in-memory \texttt{memory}
list, and \texttt{social\_neighbors}; it then initializes location against the
city map.  The core state includes affective, socioeconomic, identity,
behavioral, and information-exposure variables.  Additional mechanisms may
add namespaced or compatibility fields during setup and execution.

The class \texttt{gaworld.core.agent.Agent} is a typed adapter over this
dictionary.  It provides properties for identifier, name, state, memory,
schedule, location, and attached skill identifiers, while dictionary access
delegates to the same backing object.  \texttt{view\_as\_agent} therefore
introduces no copy, and \texttt{ensure\_agent\_dict} allows a typed caller to
hand the record back to older helpers.  This is important for incremental
migration, but it is not a fully typed agent schema: many runtime call sites
still consume dictionaries, and arbitrary mechanism-specific fields can still
be added.

Persistence extends beyond the in-memory record.  Depending on enabled
features, per-agent directories hold episodic entries, diaries, schedules,
vector-indexed recall material, environmental preferences, and other
mechanism-owned records.  On a stateful restart, stored memory is loaded into
the agent and can seed retrieval.  The file boundary makes inputs to later
decisions inspectable, although it also means that reproducibility depends on
capturing configuration, external model behavior, and the complete artifact
directory rather than only the seed CSV.

\subsection{The default per-tick path}

The cognition loop is implemented as an ordered \texttt{StagePipeline}.  A
stage receives an agent, a shared per-step dictionary, and the
\texttt{SimContext}.  Hook-visible values retain names such as
\texttt{activity}, \texttt{action}, \texttt{outcome}, and
\texttt{reflection}; intermediate values use underscore-prefixed keys.  The
default path contains the following twelve stages:

\begin{enumerate}[leftmargin=*,label=\arabic*.]
  \item \textbf{prepare} establishes tick-local inputs and emits the pre-step
  boundary;
  \item \textbf{perceive} combines local environment, social, policy, memory,
  and contributed perception sections;
  \item \textbf{interrupts} considers events that can alter the scheduled
  course;
  \item \textbf{plan} forms the agent's immediate intention in the current
  context;
  \item \textbf{adjust\_activity} reconciles that plan with scheduled and
  interrupted activity;
  \item \textbf{move} resolves travel and updates location where applicable;
  \item \textbf{select\_action} chooses an action from the available action
  space;
  \item \textbf{reflect} produces an interpretation of action and outcome;
  \item \textbf{update\_state} applies state and relationship effects;
  \item \textbf{broadcast} exposes eligible communication to other agents or
  a distributed relay;
  \item \textbf{memorize} constructs and stores the tick's experience; and
  \item \textbf{record} appends established logs, state history, and visual
  trace material.
\end{enumerate}

Before the pipeline runs, the outer loop advances the shared Clock, updates the
environment and local occupancy, polls distributed messages when enabled, and
may perform configured information seeking.  After it returns, the runtime
emits \texttt{on\_agent\_post\_step}, allowing mechanisms such as the
intervention plugin to compute their own metrics.  Day-start and day-end
events bracket the tick loop, and memory consolidation is invoked at the end
of a day.  Consequently, the twelve stages are the agent-step spine, not the
entire simulation lifecycle.

\subsection{Contribution points and compatibility calls}

The EventBus lets a mechanism observe an event, contribute zero or more items,
or filter a value in priority order.  Perception composition is a concrete
example: local context can be assembled by established code while plugins add
bounded text sections without owning the full prompt.  Post-step and day-end
events similarly let policy, economy, skills, or interests react at defined
boundaries.  Observer failures are logged and isolated unless strict mode is
enabled; stage failures instead propagate because a missing control-flow stage
could silently corrupt state.

Not all inputs use this bus.  Memory recall, environment evolution, and social
network operations retain direct paths within or around the pipeline.  The
Controller now validates structured movement requests and exposes audited
runtime interventions, but it does not mediate every selected action.
Likewise, \texttt{record} writes established artifacts in addition to the
kernel Recorder.  The lifecycle is therefore connected and configurable while
retaining compatibility paths beyond the plugin boundary.

\subsection{Configuration and ablation semantics}

The stage order is read from \texttt{CONFIG["pipeline"]["agent\_step"]}.
Entries may name built-in stages or importable callables; a custom callable can
also receive a display name.  Unresolvable entries are skipped with a warning,
and an empty resolved pipeline is rejected.  Downstream built-ins generally
read absent step values defensively, which supports mechanism ablations, but
\texttt{prepare} and \texttt{record} are intended as structural boundaries.
Changing order or removing a stage can change semantics and artifact
completeness.  The configuration surface therefore enables controlled system
experiments; it does not guarantee that every permutation is meaningful or
behavior-preserving.
